\cotefile{29}
\scanpage{95}
% Image orientation: the typescript is upside down in the supplied scan.
% archive foliation: 34
tera donc du a) dans le

\uline{Corollaire} 3.2. Soit\inserted{\unsure{en}} $S$ un préschéma non vide, $u:M'\quad M$ un homomorphisme de groupes commutatifs ordinaires, $u':G\quad G'$ l'homomorphisme transposé. Alors :
% Blank connector positions in this carbon typescript are retained as spaces.

a) Pour que $u$ soit un monomorphisme, il f et s que $u'$ soit fidèlement plat.

b) Pour que $u$ soit un épimorphisme, il f et s que $u'$ soit un monomorphisme (et alors $u'$ est même une immersion fermée).

Pour prouver a), on note que si $u$ est un monomorphisme, alors $\underline{O}_S(M)$ est un Module \deleted{\unsure{libre sur}}\inserted{sur} $\underline{O}_S(M')$, \deleted{\unsure{tout système de représentants}} admettant une base non vide (savoir, le système de sections défini par n'importe quel système de représentants dans $M$ modulo $M'$), à fortiori il est fidèlement plat. Réciproquement, s'il en est ainsi, alors $\underline{O}_S(M')\quad\underline{O}_S(M)$ est injectif, ce qui (pour $S\ne\varnothing$) implique que $u:M'\quad M$ est injectif.

Pour prouver b), on note que si $u$ est un épimorphisme, alors $\underline{O}_S(M')\quad\underline{O}_S(M)$ est surjectif, donc $u'$ est une immersion fermée et à fortiori un monomorphisme. Inversement, s'il en est ainsi, alors $\operatorname{Ker}u'={}$ groupe unité, or \deleted{\unsure{posant}} posant $M''=\operatorname{Coker}u$, on a vu que $\operatorname{Ker}u'\quad D_S(M'')$, donc par 2.1. d) on a $M''=0$ donc $u$ est un épimorphisme.

On conclut de 3.1. de la façon habituelle:

\uline{Corollaire} 3.3. Soit
\diagram{\begin{tikzcd}[column sep=small]
M' \arrow[r,phantom,"u"] & M \arrow[r,phantom,"v"] & M''
\end{tikzcd}}
une suite exacte de groupes commutatifs ordinaires, considérons la suite transposée
\diagram{\begin{tikzcd}[column sep=small]
G'' \arrow[r,phantom,"v'"] & G \arrow[r,phantom,"u'"] & G'.
\end{tikzcd}}
Alors $v'$ induit un morphisme fidèlement plat et quasi-compact de $G''$ dans $\operatorname{Ker}u'$, et ce dernier est un groupe diagonalisable.
